\begin{center}\begin{tikzpicture}[font=\small,>=stealth,line join=round]
\draw[gray!40,rounded corners=3pt] (0,0) rectangle (14.2,8.0);\node at (7.1,7.4) {两定点一样高：延长线与 x 轴永远不相交};
\begin{scope}[shift={(2.7,2.3)},x=1cm,y=1cm]\draw[->,gray] (-1.6,0)--(3.6,0) node[right]{$x$};\draw[->,gray] (0,-.35)--(0,2.8) node[left]{$y$};\draw[orange!85!black,line width=1.6pt] (0,2)--(-1,0);\draw[violet,thick] (3,2)--(-1,0);\draw[teal!80!black,thick] (0,2)--(3,2);\fill[teal!80!black] (0,2) circle(2.2pt) node[above left=4pt]{$A(0,2)$};\fill[teal!80!black] (3,2) circle(2.2pt) node[above right=4pt]{$B(3,2)$};\fill[blue] (-1,0) circle(2.2pt) node[below=4pt]{$P(-1,0)$};\end{scope}\draw[teal!80!black,dashed] (1.2,4.3)--(8,4.3);\node[anchor=west,align=left] at (8.6,4.4) {$A(0,2),\ B(3,2)$\\[12pt]$AB$ 所在直线为 $y=2$\\[12pt]与动点所在 $y=0$ 平行\\[12pt]不存在取等所需的交点。};
\node[text width=12.8cm,align=center] at (7.1,.65) {只要两定点不同且都在轴上方，距离差的两端都不能取到。};
\end{tikzpicture}\end{center}
